\section{产品设计难点：冷启动、双重冷却与真人比例}

Elys 不是一个只要模型能力足够就能自然成立的产品。它同时要解决两个冷启动：第一是用户冷启动，如何让足够多人进入网络；第二是 context 冷启动，如何让用户愿意提供足够多、足够真实、足够可用的自我信息。只有用户和 context 同时启动，赛博分身才可能带来真实连接。

\subsection{双重冷启动}

\noindent\begin{longtable}{p{0.22\textwidth}p{0.34\textwidth}p{0.35\textwidth}}
\toprule
冷启动类型 & 关键问题 & 可能解法 \\
\midrule
用户冷启动 & 没有足够真人，分身互动会显得空转 & 小圈层启动、邀请制、种子社区、场景化传播。 \\
Context 冷启动 & 用户不愿或不会交出高质量 context & 低摩擦 onboarding、即时收益、可撤回授权、示例引导。 \\
信任冷启动 & 用户担心公司掌握隐私或分身误代表自己 & 本地存储、权限分级、行为审计、敏感动作确认。 \\
网络冷启动 & 有人无关系，有内容无连接 & AI 先破冰，再引导真人确认和互动。 \\
\bottomrule
\end{longtable}

\begin{importantbox}{双重冷启动的本质}
传统社交产品主要冷启动用户和内容；AI 社交还要冷启动 context。没有 context，AI 不像你；没有真人，AI 热闹没有意义。
\end{importantbox}

\subsection{真人比例为什么重要}

Elys 的北极星指标是真人连接率或真人行为比例。这一点很关键：AI 可以降低互动门槛，但不能让网络只剩 AI。如果用户每次打开只看到 AI 回复，会产生失望，因为社交的意义仍然来自真人。AI 的角色应该是放大真人连接，而不是替代真人连接。

\begin{warningbox}{AI 社交的“热闹陷阱”}
如果产品把 AI 评论、AI 点赞、AI 发帖当成增长，短期数据可能好看，但用户会很快发现没有真人。社交网络最怕虚假繁荣，因为它会破坏信任。
\end{warningbox}

\subsection{本章小结}

Elys 的产品难点不在“能不能生成对话”，而在能否同时启动用户、context、信任和真人网络。真人连接率是比消息数更接近本质的指标。

